
\medskip
\noindent\textbf{8.4.}
\phantomsection
\label{sga3:VIIA:8.4}

In the two statements which conclude this Expos\'e, we return to
the case of a field \(k\) of arbitrary characteristic.

If \(k\) has characteristic \(0\), let \(n\ge1\). If \(k\) has
characteristic \(p>0\), let \(n\ge1\) be prime to \(p\). In both
cases we simply say that \(n\) is \emph{prime to the
characteristic of \(k\)}. If \(G\) is a group scheme over \(k\),
we denote by
\[
  n_G:G\longrightarrow G
\]
the morphism of \(k\)-schemes which sends
\(x\in G(T)\) to \(x^n\in G(T)\), for every \(k\)-scheme \(T\).

\medskip
\noindent\textbf{Proposition.}
\phantomsection
\label{sga3:VIIA:8.4-proposition}
\textit{Let \(G\) be an algebraic group over a field \(k\), and
let \(n\) be prime to the characteristic of \(k\). Then
\[
  n_G:G\longrightarrow G
\]
is an \'etale morphism.}

By VIB, \SGARefLink{prop:sga3-VIB-1.3}{1.3}, it suffices to show that \(n_G\) is \'etale at the
origin.\footnote{The editors changed ``\'etale at the origin''
to ``\'etale'' in the statement and added the preceding
sentence.} Let \(A\) be the local ring of \(G\) at the origin
and \(I\) its maximal ideal. By II, \SGARefLink{sga3:ii:corollary:3-9-4}{3.9.4}, the map
\[
  \operatorname{Lie}(n_G):
  \operatorname{Lie}(G)\longrightarrow\operatorname{Lie}(G)
\]
induced by \(n_G\) is multiplication by \(n\). It is therefore
an isomorphism, as is the endomorphism induced by \(n_G\) on
\[
  I/I^2=\operatorname{Lie}(G)^*.
\]

Suppose first that \(\operatorname{char}k=0\). The group \(G\)
is smooth over \(k\) (VIB, 1.6.1; see also VIIB, 3.3.1), and the
canonical homomorphism
\[
  \operatorname{Sym}(I/I^2)\longrightarrow
  \operatorname{gr}_I(A)
\]
is an isomorphism, where \(\operatorname{gr}_I(A)\) is the
graded ring associated with the \(I\)-adic filtration. Hence
\(n_G\) induces an automorphism of \(\operatorname{gr}_I(A)\),
and therefore also of the completion \(\widehat A\). Thus
\(n_G\) is \'etale at the origin (cf.\ EGA IV\(_4\), 17.6.3).

Suppose next that \(\operatorname{char}k=p>0\) and \(G\) has
height at most \(1\). Then \(A\) is isomorphic to the quotient
of the symmetric algebra on
\[
  \omega_{G/k}=I/I^2
\]
by the ideal generated by the \(p\)-th powers of the elements
of \(\omega_{G/k}\) (cf.\ 7.4). The same argument as in
characteristic \(0\) shows that \(n_G\) induces an automorphism
of \(A\).

Suppose that \(G\) has height at most \(r\), and that the result
has been proved for groups of height at most \(r-1\). Let
\(B,A,A'\) be the affine algebras of
\(\operatorname{Fr}G,G\), and
\[
  G'=\operatorname{Fr}G\backslash G,
\]
and let \(n_B,n_A,n_{A'}\) be the endomorphisms induced by
\(n_{\operatorname{Fr}G},n_G,n_{G'}\), respectively. Put
\[
  I'=I\cap A',
\]
the maximal ideal of \(A'\). The cartesian square
\begin{center}
  \phantomsection
  \label{sga3:VIIA:8.4-height-induction-cartesian-square}
  \begin{tikzcd}[row sep=3.4em,column sep=4.2em]
\operatorname{Fr}G \arrow[r] \arrow[d]
  & G \arrow[d] \\
e \arrow[r]
  & G'
\end{tikzcd}

\end{center}
gives
\[
  B=A/I'A.
\]
\footnote{The editors expanded the original argument in what
follows.}
The endomorphism \(n_{A'}\), respectively \(n_B\), is the
endomorphism induced by \(n_A\) on \(A'\), respectively on
\(B\). By VIA, \SGARefLink{thm:sga3-VIA-3.2}{3.2}, \(A\) is a faithfully flat
\(A'\)-module. Since \(A'\) is an Artinian local ring, because
\(G'\) is an algebraic \(k\)-group of height at most \(r-1\),
\(A\) is a free \(A'\)-module. The restriction of \(n_A\) to
\(A'\) is \(n_{A'}\), which is an isomorphism by the induction
hypothesis. Nakayama's lemma therefore shows that \(n_A\) is an
automorphism as soon as the endomorphism it induces on
\(A/I'A\) is an automorphism.

That endomorphism is \(n_B\), which is an automorphism because
\(\operatorname{Fr}G\) has height at most \(1\). Hence \(n_A\)
is an automorphism.

Finally, let \(G\) be any algebraic group over a field of
characteristic \(p>0\). The preceding argument shows that
\(n_G\) induces automorphisms of the \(k\)-schemes
\(\operatorname{Fr}^rG\). These schemes are affine over \(k\),
and their algebras are the quotients of the local algebra \(A\)
by the ideal
\[
  I^{\{p^r\}}
\]
generated by the \(p^r\)-th powers of the elements of \(I\).
Since \(n_G\) induces automorphisms of all the algebras
\(A/I^{\{p^r\}}\), passage to the projective limit shows that it
induces an automorphism of \(\widehat A\). Thus \(n_G\) is
\'etale at the origin (EGA IV\(_4\), 17.6.3).

\medskip
\noindent\textbf{8.5. Proposition.}
\phantomsection
\label{sga3:VIIA:8.5}
\textit{Let \(G\) be a finite algebraic group of rank \(n\) over
the field \(k\). Then
\[
  n_G:G\longrightarrow G
\]
is the zero morphism of \(G\).}

Combining 8.4 and 8.5 immediately gives the following
corollary.

\medskip
\noindent\textbf{Corollary 8.5.1.}
\phantomsection
\label{sga3:VIIA:8.5.1}
\textit{Let \(G\) be a finite algebraic group of rank \(n\) over
the field \(k\). If \(n\) is prime to the characteristic of
\(k\), then \(G\) is \'etale over \(k\).}
\footnote{The editors added this corollary, which the original
indicated implicitly by ``cf.\ \SGARefLink{sga3:VIIA:8.4}{8.4}.'' For another proof which
does not use \SGARefLink{sga3:VIIA:8.5}{8.5}, see \textup{[TO70]}, \SGARefLink{sga3:VIIA:5}{Lemma~5}.}

We now prove \SGARefLink{sga3:VIIA:8.5}{8.5}. Let \(H\) be a normal subgroup of \(G\), of
rank \(m\) over \(k\), and let
\[
  \lambda:H\times G\longrightarrow G
\]
be induced by the multiplication of \(G\). With the notation of
VIA, 3.2, there is a cartesian square
\begin{center}
  \phantomsection
  \label{sga3:VIIA:8.5-normal-subgroup-cartesian-square}
  \begin{tikzcd}[row sep=3.4em,column sep=4.2em]
H\times G \arrow[r,"\lambda"] \arrow[d,swap,"\operatorname{pr}_2"]
  & G \arrow[d,"\pi"] \\
G \arrow[r,"\pi"]
  & H\backslash G
\end{tikzcd}

\end{center}
Since \(\pi:G\to H\backslash G\) is faithfully flat and
quasi-compact (VIA, 3.2), while
\(\operatorname{pr}_2\) is locally free of rank \(m\), EGA
IV\(_2\), 2.5.2 shows that \(G\to H\backslash G\) is locally
free of rank \(m\). If
\[
  r=\operatorname{rk}_k(H\backslash G),
\]
then
\[
  n=\operatorname{rk}_k(G)=rm.
\]

On the other hand, for every \(k\)-scheme \(T\) there is an
exact sequence of abstract groups
\[
  1\longrightarrow H(T)\longrightarrow G(T)
  \longrightarrow(H\backslash G)(T).
\]
It is therefore clear that \(n_G\) is zero whenever \(m_H\) and
\(r_{H\backslash G}\) are zero. Taking \(H=G^0\), the identity
component, the quotient \(G^0\backslash G\) is \'etale
(cf.\ VIA, \SGARefLink{prop:sga3-VIA-5.5.1}{5.5.1}). We may consequently assume that \(G\) is
\'etale over \(k\), or that it is infinitesimal (cf.\ \SGARefLink{sga3:VIIA:7.0}{7.0}).

If \(G\) is \'etale, extension of the base field reduces us to
the case where \(k\) is algebraically closed. Then \(G\) is a
constant group (cf.\ I, \SGARefLink{sga3:i:subsection:4-1}{4.1}), and the assertion is classical.

If \(G\) is nonzero and infinitesimal, then \(k\) necessarily
has characteristic \(p>0\) (cf.\ VIB, \SGARefLink{cor:sga3-VIB-1.6.1}{1.6.1} or VIIB, \SGARefLink{sga3:VIIA:3.3.1}{3.3.1}).
The subgroups \(\operatorname{Fr}^rG\) form a composition series
of \(G\) whose successive quotients have height at most \(1\).

We are thus reduced to the case where \(G\) has height at most
\(1\). Let \(A\), respectively \(L\), be the affine algebra,
respectively the Lie algebra, of \(G\), and put
\[
  U=U_p(L).
\]
By \SGARefLink{sga3:VIIA:7.4}{7.4}, \(G=\mathfrak G_p(L)\), and hence \(A=U^*\). Therefore,
if \(\dim_kL=r\), the rank of \(G\) over \(k\) is \(p^r\)
(cf.\ \SGARefLink{sga3:VIIA:5.3.3}{5.3.3}). We shall study the morphism
\[
  p_G:G\longrightarrow G
\]
defined by taking the \(p\)-th power. It induces an
endomorphism \(p_A\) of \(A\), and, by duality, an endomorphism
\(p_U\) of \(U\).

Let \(I\) be the augmentation ideal of \(A\). We shall prove
\[
  \tag{*}\label{sga3:vii:corollary:8-5-1:equation:star}
  p_A(I)\subset I^p.
\]
Assuming this, one has
\[
  p_A^r(I)\subset I^{p^r}.
\]
Moreover,
\[
  I^{r(p-1)+1}=0,
\]
because \(I\) is generated by \(r\) elements whose \(p\)-th
powers vanish. Since \(p^r>r(p-1)\), we obtain
\(p_A^r(I)=0\), and therefore \(p_G^r\) is the zero morphism.

For every integer \(s\ge1\), let
\[
  m_A^{s-1}:A^{\otimes s}\longrightarrow A,
  \qquad
  \Delta_U^{s-1}:U\longrightarrow U^{\otimes s},
\]
denote the maps induced by the multiplication \(m_A\) of \(A\)
and the comultiplication \(\Delta_U\) of \(U\), respectively.
Then \(p_A\) is the composite
\[
  A\xrightarrow{\Delta_A^{p-1}}
  A^{\otimes p}
  \xrightarrow{m_A^{p-1}}A.
\]
Since the transpose of \(m_A\), respectively of \(\Delta_A\),
is \(\Delta_U\), respectively \(m_U\), the endomorphism
\[
  p_U={}^tp_A
\]
is the composite
\[
  U\xrightarrow{\Delta_U^{p-1}}
  U^{\otimes p}
  \xrightarrow{m_U^{p-1}}U.
\]

Let \(J\) be the augmentation ideal of \(U\).\footnote{The
editors expanded the original argument in this paragraph.}
Then
\[
  U=k1_U\oplus J,
\]
and we write \(\pi:U\to J\) for the projection with kernel
\(k1_U\). For \(s\ge1\), let \((I^s)^\perp\) be the orthogonal
of \(I^s\) in \(A^*=U\). Thus \((I^s)^\perp\) consists of the
elements \(u\in U\) for which the composite
\[
  I^{\otimes s}
  \xrightarrow{m_A^{s-1}}I
  \xrightarrow{u}k
\]
is zero. Since the transpose of \(m_A\) is \(\Delta_U\),
\((I^s)^\perp\) is the vector subspace \(P_{s-1}\) consisting
of those \(u\in U\) for which \(\Delta_U^{s-1}(u)\) vanishes on
\(I^{\otimes s}\). If
\[
  \overline\Delta_U^{s-1}
  =
  \pi^{\otimes s}\circ\Delta_U^{s-1},
\]
then
\[
  (I^s)^\perp
  =
  P_{s-1}
  =
  \ker\overline\Delta_U^{s-1}
\]
(see also VIIB, \SGARefLink{sga3:VIIB:1.3.6}{1.3.6}).

To prove \((*)\), it is enough to show that the transpose
\(p_U={}^tp_A\) maps \(P_{p-1}\) into
\[
  I^\perp=k1_U.
\]
Since \(p_U(1_U)=1_U\), it suffices to prove
\[
  \tag{**}\label{sga3:vii:corollary:8-5-1:equation:double-star}
  p_U(P_{p-1}^+)=0,
  \qquad
  P_{p-1}^+=J\cap P_{p-1}.
\]
An easy induction on \(s\) shows that \(P_{s-1}^+\) is spanned
by the products
\[
  x_1\cdots x_t,
  \qquad
  1\le t\le s-1,\quad x_i\in L
\]
(cf.\ VIIB, \SGARefLink{sga3:VIIA:4.3}{4.3}).

For \(x_1,\ldots,x_t\in L\), one has
\[
\begin{aligned}
  p_U(x_1x_2\cdots x_t)
  &=
  m_U^{p-1}
  \left(
    \prod_{j=1}^t
    \sum_{i=1}^p
    1\otimes\cdots\otimes
    \underset{\text{\scriptsize position \(i\)}}{x_j}
    \otimes\cdots\otimes1
  \right).
\end{aligned}
\]
The expression inside parentheses is a sum of \(p^t\) terms
\(x_h\), indexed by the maps
\[
  h:\{1,\ldots,t\}\longrightarrow\{1,\ldots,p\}.
\]
Such a map determines an ordered partition
\(\mathcal P_h\) of \(\{1,\ldots,t\}\) into at most \(p\)
parts.\footnote{The editors expanded the original argument and
replaced the notion of a preorder by the equivalent notion of
an ordered partition.} Indeed, let
\[
  i_1<\cdots<i_q
\]
be the elements in the image of \(h\), put
\[
  I_a=h^{-1}(i_a)
  \quad(1\le a\le q),
  \qquad
  x_{I_a}=\prod_{j\in I_a}x_j,
\]
where the product is taken in increasing order. The map \(h\)
corresponds to the \(p\)-tensor
\[
  1\otimes\cdots\otimes x_{I_1}\otimes\cdots\otimes
  x_{I_q}\otimes\cdots\otimes1,
\]
where \(x_{I_a}\) occupies position \(i_a\). Its image under the
iterated multiplication is the product
\[
  x_{I_1}\cdots x_{I_q}.
\]
This depends only on the ordered partition
\[
  \mathcal P=(I_1,\ldots,I_q);
\]
denote it by \(x_{\mathcal P}\). For a fixed \(\mathcal P\),
the same term is obtained for every choice of
\(i_1<\cdots<i_q\) in \(\{1,\ldots,p\}\). Consequently,
\[
  p_U(x_1x_2\cdots x_t)
  =
  \sum_{\mathcal P}
  \binom{p}{n(\mathcal P)}x_{\mathcal P},
\]
where \(\mathcal P\) ranges over the ordered partitions of
\(\{1,\ldots,t\}\) into at most \(p\) parts and
\(n(\mathcal P)\) is the number of parts. Here
\[
  1\le n(\mathcal P)\le\min(t,p).
\]
If \(t<p\), all the coefficients
\(\binom{p}{n(\mathcal P)}\) vanish in \(k\), and hence
\[
  p_U(x_1\cdots x_t)=0.
\]
Thus \(p_U\) vanishes on \(P_{p-1}^+\). This proves \((**)\),
then \((*)\), and completes the proof of \SGARefLink{sga3:VIIA:8.5}{8.5}.

\medskip
\noindent\textbf{Corollary 8.5.2.}\footnote{The editors added
this corollary, which is noted in \SGARefLink{sga3:viii:remark:7:3:1}{Expos\'e VIII, Remark 7.3.1}.}
\phantomsection
\label{sga3:VIIA:8.5.2}
\textit{Let \(S\) be a reduced scheme and \(G\) a finite locally
free \(S\)-group of rank \(n\). Then
\[
  n_G:G\longrightarrow G
\]
is the zero morphism of \(G\).}

Let \(S'\) be the disjoint union of the
\(\operatorname{Spec}\mathcal O_{S,\eta}\), where \(\eta\)
ranges over the maximal points of \(S\). Since \(S\) is
reduced, the morphism \(S'\to S\) is schematically dominant.
The same is therefore true of
\[
  f:G_{S'}\longrightarrow G,
\]
because \(G\) is finite locally free over \(S\) (cf.\ EGA
IV\(_3\), 11.10.5). Since \(G\to S\) is affine, and hence
separated, the coincidence locus of \(n_G\) and the zero
morphism is a closed subscheme of \(G\). By 8.5 it contains the
image of \(f\), and therefore equals \(G\). Thus \(n_G\) is the
zero morphism.

\medskip
\noindent\textbf{Remark 8.5.3.}
\phantomsection
\label{sga3:VIIA:8.5.3}

A theorem of P.~Deligne (see \textup{[TO70]}, p.~4) also says
that if \(G\) is a commutative finite locally free \(S\)-group
of rank \(n\) over an arbitrary base \(S\), then \(n_G=0\).

\chapter*{Bibliography}
\phantomsection
\label{sga3:VIIA:bibliography}
\addcontentsline{toc}{chapter}{Bibliography}

\begin{description}
  \item[\textup{[BAlg]}] N.~Bourbaki,
  \emph{Alg\`ebre}, Chapters I--III, Hermann, 1974;
  Chapter X, Masson, 1980.

  \item[\textup{[BAC]}] N.~Bourbaki,
  \emph{Alg\`ebre commutative}, Chapters I--IV, Masson, 1985.

  \item[\textup{[BLie]}] N.~Bourbaki,
  \emph{Groupes et alg\`ebres de Lie}, Chapter I,
  Hermann, 1971.

  \item[\textup{[DG70]}] M.~Demazure and P.~Gabriel,
  \emph{Groupes alg\'ebriques}, Masson \& North-Holland, 1970.

  \item[\textup{[Ja03]}] J.~C.~Jantzen,
  \emph{Representations of Algebraic Groups},
  Academic Press, 1987; 2nd ed., American Mathematical Society,
  2003.

  \item[\textup{[TO70]}] J.~Tate and F.~Oort,
  ``Group schemes of prime order,''
  \emph{Annales scientifiques de l'\'Ecole Normale
  Sup\'erieure} \textbf{3} (1970), 1--21.
\end{description}

